
\chapternonum{序言}

在“双碳”目标与数字中国建设协同推进的背景下，数字经济与能源结构优化的关系已成为经济管理与可持续发展研究中的重要议题。现有研究在计量识别方面已形成较为丰富的成果，但研究结论多以二维图表与静态文本呈现，难以直观反映省域差异、时序演化与机制路径。与此同时，增强现实技术在空间定位、三维可视化与自然交互方面持续进步，为复杂经济数据的沉浸式表达提供了新的技术路径。

本研究围绕“数字经济如何影响能源结构优化”这一核心问题，构建了“计量建模—AR可视化—多模态交互（含大语言模型）”的一体化研究框架。论文首先基于2014--2022年中国30省份面板数据，系统检验数字经济对能源结构优化的直接效应、非线性关系与中介路径；其次面向Rokid AR Studio平台设计并实现省域数据沉浸式可视化系统；在此基础上，第三章以头戴端深度感知流水线得到的手部骨骼三维坐标为基础构建手势通道，将离线语音识别与大语言模型增强相结合的语音链路与头显射线、大语言模型语义层纳入统一管线，完成指令解析、空间选择、冲突消解与智能问答及情景推演，并与设备侧空间跟踪协同支撑沉浸式分析。

本文尝试在方法层面实现经济计量研究与空间计算技术的深度融合，在应用层面提升研究结论的可读性、可解释性与可操作性。希望本研究能够为区域能源转型政策分析、沉浸式科研教学展示及相关系统开发提供可复用的思路与参考。